\begin{table}[htbp]
\centering
\caption{Swimming (moving mode) versus resting (active mode) at the baseline food/condition scenario, at the cold and warm ends of the range. $\Delta T_{\mathrm{core}}$ is the swimming minus resting defended core temperature; $\dot{Q}$ is total surface heat loss. Swimming leaves the defended core almost unchanged in all species except the smallest (minke), while raising surface heat loss by 4--28\,\%, confirming that added muscular heat is shed rather than banked --- swimming is not a thermoregulatory subsidy.}
\label{tab:swimming}
\begin{tabular}{lrrrr}
\hline
& \multicolumn{2}{c}{$T_w=-2\,^\circ$C} & \multicolumn{2}{c}{$T_w=30\,^\circ$C} \\
Species & $\Delta T_{\mathrm{core}}$ (°C) & $\Delta\dot{Q}$ (\%) & $\Delta T_{\mathrm{core}}$ (°C) & $\Delta\dot{Q}$ (\%) \\
\hline
Humpback & +0.15 & +8.5 & +0.09 & +15.0 \\
Right whale & +0.17 & +4.3 & +0.13 & +9.3 \\
Blue whale & +0.08 & +8.8 & +0.06 & +17.7 \\
Bowhead & +0.16 & +7.9 & +0.15 & +14.7 \\
Fin whale & +0.10 & +9.4 & +0.08 & +18.1 \\
Minke & +1.05 & +13.1 & +0.69 & +27.9 \\
\hline
\end{tabular}
\end{table}
